\begin{proof}
\begin{enumerate}[label=\arabic*.]
\item On the punctured base the evaluation map is
\[
 E_c:\Lambda_{\mathbb R}\longrightarrow F_c^\vee,
 \qquad E_c(x)(\omega)=\omega(x).
\]
Its kernel is zero: a real vector annihilating $F_c$ also annihilates
$\overline{F_c}$, and these two spaces span the complexified cohomology
space. Both real dimensions are six. If $\Omega$ is the evaluation
matrix of a local frame, its inverse is explicitly
\[
 E_c^{-1}(z)=R^{-1}
       \begin{pmatrix}\operatorname{Re}z\\\operatorname{Im}z\end{pmatrix},
 \qquad
 R=\begin{pmatrix}\operatorname{Re}\Omega\\\operatorname{Im}\Omega\end{pmatrix}.
\]
Under a frame change $\Omega'=G\Omega$, both $R$ and the real coordinate
column of $z$ are multiplied by the real matrix of $G$. Thus this
inverse glues. Evaluation sends each element of $\Lambda$ to its
specified period section and induces the marked real torus bundle
isomorphism, preserving the zero section. For a peripheral transport
$\gamma$ with matrix $M$, the actual evaluation square is
\[
 E_{c,\gamma p}(Mx)=\Phi_\gamma E_{c,p}(x),
\]
where $\Phi_\gamma$ is the transport in the period quotient. Consequently
\[
 J_c=E_c^{-1}iE_c=R^{-1}J_{\rm std}R,
 \qquad
 J_{\rm std}=\begin{pmatrix}0&-I_3\\I_3&0\end{pmatrix},
 \qquad J_c^2=-I_6,
\]
and $J_c(\gamma p)=M J_c(p)M^{-1}$.

\item At a finite site put $u=r e^{2\pi ia}$ on the universal angular
cover and $L=\log r/(2\pi)$. The prescribed attaching translation in
Theorem~\ref{night:P1.10} becomes
\[
 E_c^{-1}\!\left(-\frac{\log u}{2\pi i}E_c(v)\right)
       =-av+LJ_c(u)v.
\]
The affine action is $(x,u)\mapsto(Mx+v/m,\zeta_m u)$, and $Mv=v$ by
Theorem~\ref{night:P1.1}. Write
\[
 \alpha_0(x,u)=x-av,
 \qquad h_c(u)=LJ_c(u)v,
 \qquad \alpha_c=\alpha_0+h_c.
\]
The angular shift is $a\mapsto a+1/m$, and the displayed transport of
$J_c$ gives the exact identities
\[
 h_c(\zeta_m u)=M h_c(u),\qquad
 \alpha_c(Mx+v/m,\zeta_m u)=M\alpha_c(x,u).
\]
Hence $h_c$ is a section of the real local-system vector bundle on
the attaching annulus. Choose a scalar cutoff $\chi$ which is one at
the inner boundary of an attaching collar and zero at its outer
boundary. Translation by $-\chi h_c$ on that collar has inverse
translation by $\chi h_c$; the translations by $-s\chi h_c$,
$0\leq s\leq1$, give an equivariant isotopy. At the inner boundary
it carries $\alpha_c$ to $\alpha_0$ and at the outer boundary it is
the identity. The two finite collars are disjoint.

\item Write $L_k=\bigwedge^k\Lambda$ and $W_{i,k}=\bigwedge^kM_i$.
The cellular model of the main torus bundle over a graph with two
oriented edges has
\[
G_k=L_k\oplus L_{k-1}\oplus L_{k-1},\qquad
d_G(f,x_1,x_2)=((W_{1,k-1}-I)x_1+(W_{2,k-1}-I)x_2,0,0).
\tag{1.12.1}
\]
The boundary mapping torus at $i$ has
$B_{i,k}=L_k\oplus L_{k-1}$ and
\[
d_{B_i}(f,x)=((W_{i,k-1}-I)x,0).
\tag{1.12.2}
\]
The finite boundary maps to $G$ by identity on $L_k$ and inclusion
of the suspension into its respective edge.  The cusp boundary,
corresponding to $(g_1g_2)^{-1}$, maps by
\[
j_{0,G}(f,x)=(f,-W_{2,k-1}W_{0,k-1}x,-W_{0,k-1}x).
\tag{1.12.3}
\]
This follows by traversing the second edge backwards and then the
first backwards, transporting the fibre chain to each initial end.
Indeed its boundary is $W_{0,k-1}-I$ because
\[
-(W_1-I)W_2W_0-(W_2-I)W_0=W_0-I.
\]

\item Here are full models for the finite pieces and their attaching maps.
Put $\varepsilon_i=\ell(v_i)$ and
$P_i=(\varepsilon_iv_i,e_1,\ldots,e_5)$.
This is unimodular.  It splits
$\Lambda=\mathbb Z(\varepsilon_iv_i)\oplus K$, $K=\ker\ell$,
and $P_i^{-1}M_iP_i=1\oplus R_i$.
The cap retracts onto the mapping torus of $R_i$ on $T_K$.
Its chains are
\[
E_{i,k}=\bigwedge^kK\oplus\bigwedge^{k-1}K,
\qquad d_{E_i}(f,x)=((\bigwedge^{k-1}R_i-I)x,0).
\tag{1.12.4}
\]
In these split coordinates, with $v=\varepsilon_iv_i$, the finite
boundary-to-cap chain map is specified on all summands by
\[
\begin{array}{c|c}
\text{boundary chain}&\text{cap chain}\\\hline
x\in\bigwedge^kK&(x,0)\\
v\wedge x&(0,\varepsilon_i\sum_{j=0}^{m_i-1}
                  (\bigwedge^{k-1}R_i)^jx)\\
g_i\wedge x&(0,x)\\
g_i\wedge v\wedge x&0.
\end{array}
\tag{1.12.5}
\]
Changing back uses
$\operatorname{diag}(\bigwedge^kP_i^{-1},
\bigwedge^{k-1}P_i^{-1})$.
Geometrically this map is
$(t,y,\theta)\mapsto[y,\theta+\varepsilon_im_it]$ into the mapping
torus.  The $t$-circle sweeps $m_i$ strips, giving the norm sum; the
$g_i$-circle sweeps one strip; their product maps to a single
base direction, giving zero on the last row.  This proves the map
including its sign and the possible finite-cap torsion.

\item We justify that these are integral cellular models, rather than a
formal use of invariant rational forms.  In the bases with columns
\[
K_1=\begin{pmatrix}
-1&-1&0&-1&0\\0&1&0&1&1\\0&1&0&0&0\\
0&-1&-1&0&-1\\0&0&1&0&0
\end{pmatrix},\qquad
K_2=\begin{pmatrix}
-1&0&-1&0&-1\\-1&1&0&1&-1\\0&1&-1&1&-1\\
0&-1&0&-1&1\\1&-1&1&0&1
\end{pmatrix},
\]
both unimodular, $R_1$ is the permutation of five coordinates in a
cycle, and $R_2$ is a four-cycle together with multiplication by $-1$
on the fifth coordinate.  The minimal product cell structure of
circles is preserved by these signed permutations; its cellular
boundary is zero and its action on cells is exactly the exterior
power.  Formula (1.12.5) is obtained by subdividing its swept
base interval into the indicated strips.  The common torus can be
marked into these two structures by based cellular approximations.
Their cellular maps are the exterior-power matrices because the
cellular differentials vanish.  Based maps between these aspherical
tori with a fixed map on $\pi_1$ have a contractible space of
homotopies: lifting to the universal covers gives equivariant maps
to a contractible vector space and linear interpolation relative to
the base point.  Thus the choices introduce no extra attaching class
on the graph.  Composing the two edge cylinders gives exactly
(1.12.3).  The explicit conjugators were found in the bounded search
$[-1,1]^5$ and then verified by exact multiplication in
\texttt{finite\_permutation\_models.py}.

\item Let $F_k:B_{1,k}\oplus B_{2,k}\to G_k\oplus E_{1,k}\oplus E_{2,k}$
have the main boundary maps as its first row, and the negatives of
(1.12.5) in the corresponding two cap rows.  The chain complex of
$Y$ is the mapping cone
\[
C_k(Y)=G_k\oplus E_{1,k}\oplus E_{2,k}
       \oplus B_{1,k-1}\oplus B_{2,k-1},\qquad
d_Y=\begin{pmatrix}d_G\oplus d_{E_1}\oplus d_{E_2}&F_{k-1}\\
0&-d_{B_1}\oplus-d_{B_2}\end{pmatrix}.
\tag{1.12.6}
\]
Equations (1.12.1)--(1.12.6), minors for exterior powers, and the input
matrices specify every entry of every differential.
The full arrays are exported to \texttt{complement\_export.json}; its
program verifies $d^2=0$ in all degrees and computes Smith
decompositions.  Chain ranks in degrees zero through eight are
$(3,22,71,132,155,118,57,16,2)$.
The resulting homology of $Y$ is free, with ranks
$(1,1,3,9,12,9,3,0,0)$.

\item The boundary $B_0$ has
\[
H_1(B_0)=\mathbb Z^4,\qquad H_2(B_0)=\mathbb Z^{12},\qquad
H_3(B_0)=\mathbb Z^{18}\oplus\mathbb Z/2.
\]
The homomorphism $j_Y:H_k(B_0)\to H_k(Y)$ from (1.12.3) is
surjective for $k=1,2,3$.  In degrees two and three, the program
\texttt{complement\_kernel\_check.py} supplies integral cycle bases
of its kernel whose entire suspension block in (1.12.2) is zero.
This is an exact zero matrix, not a vanishing after tensoring with
$\mathbb Q$.  Therefore these kernels are represented by cycles of
the nearby torus itself.  The restriction of the actual map to the
cusp cap on these kernels is consequently the already proved
collapse of P1.11, regardless of its value on other boundary
suspension cycles.

\item For transparency, in the Smith bases exported with these cycles the
combined map in degree one is
\[
\begin{pmatrix}20&1&0&0\\1&0&-1&-1\\1&0&0&0\\0&0&2&1\end{pmatrix},
\qquad\det=1.
\tag{1.12.7}
\]
The cusp meridian maps to zero here because the zero section extends.
In degree two $\ker j_Y=\mathbb Z^9$ and its map to $H_2(C)$ is
\[
\begin{pmatrix}
0&0&1&-1&0&-1&0&0&0\\
0&0&1&0&0&-1&0&1&0\\
0&0&1&-1&0&-2&0&1&0\\
0&0&0&0&-1&-1&0&0&-1\\
1&1&-1&0&0&1&1&0&0\\
1&2&0&0&1&0&0&0&0\\
-1&-1&0&0&0&0&0&0&0\\
0&0&0&0&1&0&0&0&0\\
0&0&0&0&0&-1&0&0&0
\end{pmatrix},\qquad\det=1.
\tag{1.12.8}
\]
In degree three $\ker j_Y=\mathbb Z/2\oplus\mathbb Z^9$.
Its collapse matrix, with the torsion generator first, is
\[
\begin{pmatrix}
0&1&0&1&0&0&0&0&0&1\\
0&1&-1&0&0&0&-1&0&-1&0\\
0&0&0&0&0&0&1&0&1&0\\
0&0&0&-1&1&-1&0&-1&0&0\\
0&0&0&0&0&0&0&1&0&0\\
0&0&0&0&0&0&1&0&0&0\\
0&0&0&0&0&1&0&1&0&0\\
0&0&1&0&1&-2&1&-1&2&0\\
0&1&0&-1&1&-2&0&-1&0&0
\end{pmatrix}.
\tag{1.12.9}
\]
The determinant of its last nine columns is one.  In particular the
torsion of the boundary and the second finite cap has been included
throughout the calculation.

\item We now apply Mayer--Vietoris to $M_c=Y\cup C_{\rm cap}$.
Because $j_Y$ is onto, subtracting a preimage of the $Y$-coordinate
identifies the cokernel of $(j_Y,j_C)$ with
$H_k(C)/j_C(\ker j_Y)$; its kernel is
$\ker j_Y\cap\ker j_C$.
Equations (1.12.7)--(1.12.9) therefore say that the combined map is an
isomorphism in degrees one and two and is onto in degree three.
The exact sequence yields $H_2(M_c)=H_3(M_c)=0$.
We already proved $\pi_1(M_c)=0$, and successive applications of
Hurewicz now give $\pi_2(M_c)=\pi_3(M_c)=0$.

\item Integral Poincar\'e duality and the universal coefficient theorem
give $H_5=H_6=H_7=0$.  They also show that $H_4$ is free: its possible
torsion would be $\operatorname{Ext}(H_4,\mathbb Z)$ in $H^5$, which
vanishes by duality with $H_3=0$.
Finally $\chi=6$ gives $\operatorname{rank}H_4=4$.
This proves all claimed integral homology groups without making an
assumption about the higher-degree cusp suspension chain term.

\end{enumerate}
\end{proof}
